\chapter{Performance Optimized Organization: Pipelining and Instruction-Level Parallelism}\label{lect5}
\index{pipelining!performance optimized organization}\index{ILP (Instruction-Level Parallelism): }\index{performance optimized organization}
\localtoc

In Chapter \ref{lect4}, we examined component-level digital logic optimizations---such as carry-lookahead adders, Kogge-Stone parallel prefix networks, and Wallace tree multipliers---designed to minimize combinational propagation delay along critical paths within individual functional blocks. In this chapter, we elevate our perspective from circuit-level logic to \textbf{system- and microarchitectural-level organization}. Rather than merely accelerating an isolated combinational operation, performance-optimized organization restructures the cooperative flow of instructions through time and space.

Instruction-Level Parallelism (ILP) represents the family of architectural techniques that enable multiple machine operations to be executed simultaneously. We focus on the foundational organizational mechanisms that govern modern processor throughput: pipelining, pipeline hazard management (structural, data, and control), dynamic branch prediction, and superscalar dispatch with dynamic scheduling. As we will discover, every performance gain realized through spatial or temporal overlap introduces architectural dependencies that demand rigorous hardware solutions.

\section{Pipelining}

Combinational logic circuits (CLCs) that have a very high depth cause the system clock frequency to be reduced. The solution that allows for an increase in clock frequency is the introduction of pipeline registers.


\subsection{Pipeline Acceleration\index{pipeline acceleration}}

In Figure \ref{fig47} the pipeline technique is illustrated. In Figure \ref{fig47}a, the frequency at which the system clock can work is given by the propagation time through the $CLC_{f \circ g}$ to which is added the time associated with the propagation through the registers.

\placedrawing[htbp]{figures/05b_01_pipelining.lp}{Pipelining. {\bf a.} The initial circuit, without pipeline register.  {\bf b.} The pipelined version of the circuit.}{fig47}

\begin{equation}
    T_{clock} = t_{regProp} + t_{CLC_{f \circ g}} + t_{set-up} \simeq t_{CLC_{f \circ g}}
\end{equation}

In Figure \ref{fig47}b, the frequency at which the system clock can work will be increased because the period of the clock is given by:
\begin{equation}
    T_{clock} = t_{regProp} + \max(t_{CLC_{f}}, t_{CLC_{g}}) + t_{set-up}
\end{equation}

The resulting acceleration (speedup) is:
\begin{equation}
    \alpha \simeq \frac{t_{CLC_{f \circ g}}}{\max(t_{CLC_{f}}, t_{CLC_{g}})}
\end{equation}

The maximum efficiency, $\alpha \simeq 2$, is obtained when the combinational delay is perfectly balanced across stages:
\begin{equation}
    t_{CLC_{f}} \simeq t_{CLC_{g}}
\end{equation}

If the resulting frequency is not high enough, the technique is applied recursively to the deepest circuit or both. The process can continue until the requested speed is reached or until register setup and propagation overheads dominate the stage delay.

\subsection{Pipeline Latency}

Increasing the clock frequency is accompanied, in a pipeline structure, by the introduction of additional latency. By latency, we mean the number of clock cycles that occur in the sequential processing of a data string.

For example, in the structure from Figure \ref{fig47}a the latency is 2, due to the two synchronization registers, on the input and output. By introducing a pipeline register in the structure in Figure \ref{fig47}b the latency increases by one unit.

If we process a sufficiently long string of data, \texttt{X(1), X(2), ..., X(n)}, then the additional latency introduces an insignificant effect: the processing of the string will be performed in $n+3$ cycles instead of $n+2$ but at a frequency $\alpha$ times higher.

In other words, the introduction of the pipeline register in Figure \ref{fig47}b triggers a temporal parallel process wherein the system concurrently operates on two successive components of the data sequence.

\subsection{Pipeline Throughput and Performance Metrics}
\index{CPI (Cycles Per Instruction)!pipeline performance}\index{pipeline speedup}\index{performance equation}

The performance of an instruction pipeline is formally evaluated using the fundamental processor performance equation:
\begin{equation}
    \text{Execution Time} = \text{IC} \times \text{CPI} \times T_{clock}
\end{equation}
where $\text{IC}$ is the total instruction count of the program, $\text{CPI}$ is the average number of clock cycles per instruction, and $T_{clock} = 1/f_{clock}$ is the clock period.

In an ideal $k$-stage pipeline operating under steady-state conditions without hazards or stalls, exactly one instruction completes on every clock cycle:
\begin{equation}
    \text{CPI}_{\text{ideal}} = 1
\end{equation}
However, real hardware pipelines deviate from this theoretical optimum due to pipeline stalls introduced by architectural dependencies:
\begin{equation}
    \text{CPI} = \text{CPI}_{\text{ideal}} + \text{Stalls}_{\text{structural}} + \text{Stalls}_{\text{data}} + \text{Stalls}_{\text{control}}
\end{equation}

The effective speedup $S_k$ of a $k$-stage pipeline over an unpipelined single-cycle implementation is given by:
\begin{equation}
    S_k = \frac{\text{Execution Time}_{\text{unpipelined}}}{\text{Execution Time}_{\text{pipelined}}} = \frac{\text{CPI}_{\text{unpipelined}} \times T_{\text{unpipelined}}}{\text{CPI}_{\text{pipelined}} \times T_{\text{pipelined}}} \simeq \frac{k}{1 + \text{Stalls}}
\end{equation}
Minimizing these pipeline stall penalties through performance-optimized organization is the central engineering objective of high-performance processor design.


\section{{\bf {\em Pipelined Version of toyRISC}}}

\subsection{Structure}

The operating frequency of the toyRISC processor presented section \ref{toyrisc} can be increased by inserting some pipeline registers on the critical path of the combinatorial propagation. A version, adapted to be accelerated, is shown in Figure \ref{fig46} where the two memories, for data and programs, are also represented. The program memory is a synchronous memory (see \ref{syncMem}) connected in our system  combinationally  through the output multiplexers. The data memory is a synchronous pipelined memory (see \ref{sybcPipeMem}) which responds with one cycle latency to the read command.

The micro-architecture is described is the same as the one presented in Listing \ref{lst:toyRISCdefines} from the chapter \ref{lect4}. The same for the Instruction Set Architecture which is in Table \ref{tab:riscISA}.

\placedrawing[htbp]{figures/05b_02_pipelined_toyRISC.lp}{The pipelined version of toyRISC.}{fig46}

The three pipeline registers introduced in design divide the execution of an instruction in four stages:
\begin{description}
    \item[INSTRUCTION FETCH (IF)]: the content of the register {\tt pc} (program counter) addresses in the program memory the binary code of an instruction; in the same time the module {\tt next} computes the value for the next {\tt pc}. The execution time for this stage is:
        $$t_{IF} = t_{pc} + max(t_{next}, t_{ACCto Program Memory}) + t_{suPipeReg1}$$
        In the pipeline register {\bf PipeReg\_1} is loaded the information requested to finish the execution of the fetched instruction
        
    \item[OPERANDS FETCH (OF)]: from the register are fetched the two operands of the instruction using the fields {\tt leftAddr} and {\tt rightAddr} fetched in the previous cycle from the program memory. In {\bf pipeReg\_2} are loaded the two fetched operands and the information needed in the next cycles to end the execution of the instruction. The execution time for this stage is:
        $$t_{OF} = t_{PipeReg1} + t_{regFile} + t_{suPipeReg2}$$


        
    \item[EXECUTION (EX)]: ALU performs the operation according to the code {\tt opCode}, propagated to this stage through the two previous  pipeline registers,  generate {\tt result} and the predicate {\tt zero = (result == 0)}; the {\tt add} module adds to the program counter the signed value of {\tt value} to perform the branch operation.  The execution time for this stage is:
        $$t_{EX} = t_{PipeReg2} + max((t_{mux}+t_{alu}),  (t_{add}+t_{mux}), t_{dataMemPipeReg}) + t_{suPipeReg3}$$
    \item[WRITE BACK (WB)]: acts on the first two levels in the pipeline: it allows the modification of the {\tt pc} to perform jumps and writes the result of the current operation in the register file. The execution time for this stage is:
        $$t_{WB} = t_{PipeReg3} +  max(t_{suRegFile}, (t_{next} + t_{suPC}))$$
\end{description}
The maximum clock frequency is limited by:
$$T_{clock} = max(t_{IF}, t_{OF}, t_{EX}, t_{WB})$$



\subsection{Latency\index{latency}}

In each clock cycle the execution of one instruction ends with a latency of three clock cycles. The entire structure performs in parallel 4 successive instructions, each stage being involved in the execution of one instruction.

\begin{exemplu}
Let us see how is executed the following sequence of instructions:
\begin{verbatim}
    XOR(5,0,0)  ;
    VAL(1,12)   ;
    SUB(2,3,4)  ;
    ADD(0,0,3)  ;
    AND(1,3,2)  ;
\end{verbatim}
In Table \ref{pipeEx} the distribution along the pipe of the execution is represented. Each line in the table represents a pipeline level. The execution of the first instruction, {\tt XOR}, starts in the $i$-th clock cycle, with {\tt XOR0}, and ends in the $(i+3)$-th cycle with {\tt XOR3}. Then in each clock cycle one of the following instruction ends.

\begin{table}[H]
  \begin{center}
    \caption{Pipelined execution.}
    \label{pipeEx}
    {\small
    {\tt
    \begin{tabular}{|c||c|c|c|c|c|c|c|c|}
      \hline
     Pipe Stage $\setminus$ Time    & t=i & t=i+1 & t=i+2 & t=i+3 & t=i+4 & t=i+5 & t=i+6 & t=i+7   \\ \hline \hline
        Pipe0:IF        & XOR0 & VAL0   & SUB0   & ADD0   & AND0   &      &       &         \\ \hline
        Pipe1:OF        &     & XOR1   & VAL 1  & SUB1   & ADD1   & AND1   &          &         \\ \hline
        Pipe2:EX        &     &       & XOR2   & VAL2   & SUB2   & ADD2   & AND2   &        \\ \hline
        Pipe3:WB        &     &       &       & XOR3   & VAL3   & SUB3   & ADD3   & AND3        \\ \hline
    \end{tabular}
    }
    }
  \end{center}
\end{table}
$\diamond$
\end{exemplu}

In each clock cycle, an instruction completes with a latency of 3 clock cycles. The good news: the clock frequency increases significantly due to pipeline organization. The bad news: the latency introduced by the pipeline organization creates unwanted dependencies in single-threaded execution.

\textbf{Architectural Insight: Hazard-Free Execution via Barrel Multithreading}:
It is instructive to note a fundamental architectural relationship between pipeline depth and multithreading. Because this pipeline has exactly $k = 4$ stages, an instruction completes its write-back step exactly 4 clock cycles after it is fetched. In a single-threaded program, an instruction immediately following a producer will encounter data or control hazards. However, if the processor core interleaves $N = 4$ independent hardware threads in a fixed round-robin order (known as \textit{barrel} or \textit{interleaved multithreading}), each individual thread issues an instruction only once every 4 cycles. As a result, when thread $T_0$ issues its second instruction at cycle $t=4$, the result of its first instruction (issued at $t=0$) has already been written back to the register file in stage WB at cycle $t=3$. Under this condition---where the number of interleaved threads equals the pipeline depth ($N = k$)---all pipeline data and control hazards are structurally eliminated without needing stall bubbles, interlock logic, or forwarding multiplexers. We will explore this architectural technique in detail in later chapters.

\section{Hazards Generated by Dependencies}

There are three fundamental types of dependencies that create pipeline hazards\index{pipeline dependencies: structural, data, and control hazards}:
\begin{enumerate}
\item \textbf{Structural Hazard}: Hardware resource conflicts where multiple instructions attempt to use the same physical unit concurrently.
\item \textbf{Data Hazard}: Dependencies where an instruction requires the result of an earlier instruction that has not yet completed write-back.
\item \textbf{Control Hazard}: Dependencies where the next instruction address depends on the evaluation of a branch or jump in a subsequent stage.
\end{enumerate}

These dependencies generate hazardous behaviors in the pipelined processor due to the temporal overlap of partially completed instructions.

\subsection{Structural Hazards and Hardware Replication}\label{sec:structural_hazards}
\index{structural hazards!resource conflicts}\index{hardware replication!structural hazard mitigation}

A \textbf{structural hazard} occurs when two or more overlapping instructions in the pipeline attempt to access the same physical hardware resource simultaneously. When a required hardware block cannot support all concurrent requests, one or more instructions must stall until the resource is freed, degrading processor throughput.

\subsubsection{Pipelined Structural Hazards and Dedicated PC Adder Replication}
In a non-pipelined, single-cycle computing system (such as the Harvard 5-OS core discussed in Chapter \ref{lect6a}), execution is strictly sequential. Consequently, a single arithmetic logic unit (ALU) can theoretically be time-multiplexed to perform both arithmetic computations and sequential program counter incrementing ($\text{pc} \leftarrow \text{pc} + 4$).

In a pipelined processor, however, this time-multiplexing is impossible. On \textbf{every single clock cycle}, the Instruction Fetch (\texttt{IF}) stage must calculate the sequential next address ($\text{pc} + 4$) concurrently while the Execution (\texttt{EX}) stage evaluates an arithmetic instruction or computes a branch target ($\text{pc} + \text{offset}$). If the processor attempted to share a single ALU between \texttt{IF} and \texttt{EX}, a structural hazard would occur on every cycle, stalling instruction fetch whenever arithmetic was performed and halving throughput.

To eliminate this structural hazard, pipelined organization mandates \textbf{hardware replication}:
\begin{itemize}
    \item \textbf{Dedicated Program Counter Adder}: The \texttt{IF} stage incorporates its own autonomous $+4$ incrementer (or $+1$ in word-addressed architectures) exclusively dedicated to sequential instruction pointer calculation.
    \item \textbf{Dedicated Branch Target Adder}: A separate adder is provided in the \texttt{EX} (or \texttt{OF}) stage to compute relative branch offsets without contending with the primary integer ALU.
    \item \textbf{Separate Instruction and Data Memories}: In von Neumann architectures with a single shared memory bus, a structural hazard arises whenever an instruction in \texttt{MEM} attempts to read/write data while \texttt{IF} attempts to fetch the next instruction. By utilizing the Harvard architecture with separate instruction and data memories (or dual-ported L1 caches), memory access conflicts are eliminated.
    \item \textbf{Multi-Ported Register Files}: The register file must provide multiple independent read ports (at least two for binary operations in \texttt{OF}) and a concurrent write port in \texttt{WB} so that operand reading and result write-back do not create structural contention.
\end{itemize}

\subsubsection{Superscalar Structural Hazards and Execution Unit Replication}
When the microarchitecture is expanded from a scalar pipeline (issuing 1 instruction per cycle) to a \textbf{superscalar} organization (issuing $N$ instructions per cycle), structural hazards multiply:
\begin{itemize}
    \item If two independent arithmetic instructions are issued simultaneously in the same cycle, a core with only one integer ALU suffers a structural stall.
    \item If two memory operations (e.g., two loads, or a load and a store) are issued in parallel, they conflict unless the data cache provides multiple independent read/write ports.
\end{itemize}
To sustain high issue rates without structural stalls, superscalar processors replicate functional execution units:
\begin{itemize}
    \item Multiple symmetric integer ALUs (e.g., 2 to 4 ALUs operating in parallel).
    \item Multiple Address Generation Units (AGUs) and dual-ported load/store queues.
    \item Dedicated floating-point adders and multipliers separate from the integer datapath.
    \item Highly multi-ported register files (requiring $2N$ read ports and $N$ write ports for an $N$-way superscalar core).
\end{itemize}
Where physical silicon area or routing complexity prevents full replication (e.g., hardware dividers or complex transcendental units), instructions targeting the shared unit must arbitrate, and subsequent instructions that require the busy unit experience a structural stall.

\subsection{Data dependency\index{dependencies!data dependencies in pipelined processors}}

When the current instruction uses values generated by a previous instruction that has not been completed, the effect of data dependency appears. There are 3 types of data dependencies that
we’ve been talking about:
\begin{itemize}
\item RAW: Read after Write
\item WAR: Write after Read
\item WAW: Write after Write
\end{itemize}
We will focus on the first, RAW, which is typical for our processor version. For the other two, they are illustrated in Example \ref{warwaw}.

\begin{exemplu}\label{dataDep}
Let us see how is executed the following sequence of instructions:
\begin{verbatim}
    XOR(5,0,0)  ;
    VAL(1,12)   ;
    SUB(2,3,4)  ;
    ADD(0,2,1)  ;
    AND(1,1,2)  ;
\end{verbatim}
In Table \ref{pipeEx2} the distribution along the pipe of the execution is represented.

\begin{table}[H]
  \begin{center}
    \caption{Data dependency in the pipelined execution .}
    \label{pipeEx2}
    {\small
    {\tt
    \begin{tabular}{|c||c|c|c|c|c|c|c|c|}
      \hline
     Pipe Stage $\setminus$ Time  & t=i   & t=i+1 & t=i+2 & t=i+3  & t=i+4  & t=i+5  & t=i+6 & t=i+7    \\ \hline \hline
        Pipe0:IF                  & XOR0  & VAL0  & SUB0  & ADD0  & AND0   &        &        &                \\ \hline
        Pipe1:OF                  &       & XOR1  & VAL1  & SUB1  & ADD1   & AND1   &        &                 \\ \hline
        Pipe2:EX                  &       &       & XOR2  & VAL2  & SUB2   & {\footnotesize {\bf {\em ADD2}}}   & AND2   &                 \\ \hline
        Pipe3:WB                  &       &       &       & XOR3  & VAL3   & {\footnotesize {\bf {\em SUB3}}}   & ADD3   & AND3              \\ \hline
    \end{tabular}
    }
    }
  \end{center}
\end{table}
The WB cycle for SUB, {\tt SUB3}, is executed at {\tt i+5}, too late to have the content of {\tt rf2} actualized with the result of  {\tt SUB(2,3,4)} for the instruction {\tt ADD(0,2,1)} which is supposed to have {\tt rf2} actualized by the previous instruction in {\tt i+4}. {\tt ADD2} must be preceded by {\tt SUB3}. {\tt SUB3} is delayed with one clock cycle for a proper execution of the code. The instruction {\tt ADD(0,2,1)} uses the value in {\tt rf2} before the execution of {\tt SUB{2,3,4}}.

$\diamond$
\end{exemplu}


\begin{exemplu}\label{dataDepSim}
Data dependency is illustrated also by running on our simulator the  program run in Example \ref{exADD}:
\begin{verbatim}
            NOP         ;
            VAL(0,1)    ;
            VAL(1,2)    ;
            VAL(2,3)    ;
            VAL(3,4)    ;
            VAL(4,5)    ;
            NOP         ;
            NOP         ;
            ADD(5,4,3)  ;
            HALT        ;
            HALT        ;
\end{verbatim}
The result is correct because of the two {\tt NOP}s inserted in the program from Example \ref{exADD} before the {\tt ADD} instruction:
{\small
\begin{verbatim}
pc=  0  RF=[x, x, x, x, x, x, x, x] leftOp2=x rightOp2=x result3=x
pc=  1  RF=[x, x, x, x, x, x, x, x] leftOp2=x rightOp2=x result3=x
pc=  2  RF=[x, x, x, x, x, x, x, x] leftOp2=x rightOp2=x result3=x
pc=  3  RF=[x, x, x, x, x, x, x, x] leftOp2=x rightOp2=x result3=x
pc=  4  RF=[x, x, x, x, x, x, x, x] leftOp2=x rightOp2=x result3=1
pc=  5  RF=[1, x, x, x, x, x, x, x] leftOp2=x rightOp2=x result3=2
pc=  6  RF=[1, 2, x, x, x, x, x, x] leftOp2=1 rightOp2=1 result3=3
pc=  7  RF=[1, 2, 3, x, x, x, x, x] leftOp2=1 rightOp2=1 result3=4
pc=  8  RF=[1, 2, 3, 4, x, x, x, x] leftOp2=1 rightOp2=1 result3=5
pc=  9  RF=[1, 2, 3, 4, 5, x, x, x] leftOp2=1 rightOp2=1 result3=1
pc= 10  RF=[1, 2, 3, 4, 5, x, x, x] leftOp2=5 rightOp2=4 result3=1
pc= 10  RF=[1, 2, 3, 4, 5, x, x, x] leftOp2=1 rightOp2=1 result3=9
pc= 10  RF=[1, 2, 3, 4, 5, 9, x, x] leftOp2=1 rightOp2=1 result3=1
\end{verbatim}
}
If the two {\tt NOP}s are omitted,
\begin{verbatim}
            NOP         ;
            VAL(0,1)    ;
            VAL(1,2)    ;
            VAL(2,3)    ;
            VAL(3,4)    ;
            VAL(4,5)    ;
            ADD(5,4,3)  ;
            HALT        ;
            HALT        ;
\end{verbatim}
then the processor behaves incorrectly, as follows:
\begin{verbatim}
pc=   0  RF=[x, x, x, x, x, x, x, x] leftOp2=x rightOp2=x result3=x
pc=   1  RF=[x, x, x, x, x, x, x, x] leftOp2=x rightOp2=x result3=x
pc=   2  RF=[x, x, x, x, x, x, x, x] leftOp2=x rightOp2=x result3=x
pc=   3  RF=[x, x, x, x, x, x, x, x] leftOp2=x rightOp2=x result3=x
pc=   4  RF=[x, x, x, x, x, x, x, x] leftOp2=x rightOp2=x result3=1
pc=   5  RF=[1, x, x, x, x, x, x, x] leftOp2=x rightOp2=x result3=2
pc=   6  RF=[1, 2, x, x, x, x, x, x] leftOp2=1 rightOp2=1 result3=3
pc=   7  RF=[1, 2, 3, x, x, x, x, x] leftOp2=1 rightOp2=1 result3=4
pc=   8  RF=[1, 2, 3, 4, x, x, x, x] leftOp2=x rightOp2=x result3=5
pc=   8  RF=[1, 2, 3, 4, 5, x, x, x] leftOp2=1 rightOp2=1 result3=x
pc=   8  RF=[1, 2, 3, 4, 5, x, x, x] leftOp2=1 rightOp2=1 result3=1
\end{verbatim}
because the operands for the arithmetic operations are not yet in place being delayed on the execution pipe.

$\diamond$
\end{exemplu}



To solve the problem of data dependency just emphasized there are used two solutions: stalling and forwarding.

\subsubsection{Stalling\index{stalling: solution for pipelined processor dependency}}

An inefficient solution, for the execution time, is to stall by introducing a {\tt NOP} instruction between {\tt SUB(2,3,4)} and {\tt ADD(0,2,1)}, thus delaying the use of the content of {\tt rf2} with one clock cycle. Result the following sequence of instructions:
\begin{verbatim}
    XOR(5,0,0)  ;
    VAL(1,12)   ;
    SUB(2,3,4)  ;
    NOP         ;
    ADD(0,2,1)  ;
    AND(1,1,2)  ;
\end{verbatim}
whose execution is illustrated in Table \ref{stall}.

\begin{table}[H]
  \begin{center}
    \caption{Stalling}
    \label{stall}
    {\small
    {\tt
    \begin{tabular}{|c||c|c|c|c|c|c|c|c|c|}
      \hline
     Instr $\setminus$ Time & t=i & t=i+1 & t=i+2 & t=i+3 & t=i+4 & t=i+5 & t=i+6 & t=i+7 & t=t+8   \\ \hline \hline
        IF      & XOR0 & VAL0   & SUB0   & NOP0   & ADD0   & AND0   &     &       &         \\ \hline
        OF      &     & XOR1   & VAL1   & SUB1   & NOP1   & ADD1   & AND1   &         &         \\ \hline
        EX      &     &       & XOR2   & VAL2   & SUB2   & NOP2   & {\footnotesize {\bf {\em ADD2}}}   & AND2   &       \\ \hline
        WB      &     &       &       & XOR3   & VAL3   & {\footnotesize {\bf {\em SUB3}}}   & NOP3   & ADD3   & AND3       \\ \hline
    \end{tabular}
    }
    }
  \end{center}
\end{table}
Now, {\tt ADD2} is preceded by {\tt SUB3} and addition is performed takeing into account the result of the subtract operation.

\subsubsection{Reordering\index{reordering: solution for pipelined processor dependency}}

Sometimes, {\bf but only sometimes}, there is a very simple and efficient solution for data dependency: reordering the instructions in the sequence of instructions. In Example \ref{dataDep}, the instruction {\tt XOR(5,0,0)} can be moved, without affecting the result of running the sequence of instructions, between the instructions {\tt SUB} and {\tt ADD}. as follows:
\begin{verbatim}
    VAL(1,12)   ;
    SUB(2,3,4)  ;
    XOR(5,0,0)  ;
    ADD(0,2,1)  ;
    AND(1,1,2)  ;
\end{verbatim}
Thus, the {\tt XOR} instruction is used for stalling instead of the {\tt NOP} instruction, but now the execution time is not affected.

\subsubsection{Forwarding\index{forwarding: solution for pipelined processor dependency}}

By adding hardware, the data dependency can be resolved, {\bf in all cases}, with no time penalty. The solution is to connect {\tt result} not only to the register file input, but also directly to the ALU input when the value just computed is needed for the next instruction. The mechanism is called {\em forwarding} and requires the addition of one input to the  ALU input as left operand or as right operand. But, because there are binary operations (operations with two operands) sometimes both operands arrive to late in the register file to be fetched for the current instruction.

Thus, forwarding control is done by a circuit which analyses the binary code of three successive instructions. Three situations can occur:

\begin{verbatim}
    xxxxxx_00001_xxx...x                // identified in PipeReg_3
    xxxxxx_yyyyy_00001_xxx...x          // identified in PipeReg_2
\end{verbatim}
or
\begin{verbatim}
    xxxxxx_00001_xxx...x                // identified in PipeReg_3
    xxxxxx_yyyyy_zzzzz_00001_xxx...x    // identified in PipeReg_2
\end{verbatim}
or
\begin{verbatim}
    xxxxxx_00001_xxx...x                // identified in PipeReg_3
    xxxxxx_yyyyy_00001_00001_xxx...x    // identified in PipeReg_2
\end{verbatim}

\begin{verbatim}
    xxxxxx_00001_xxx...x                // identified in PipeReg_4
    xxxxxx_00011_xxx...x                // identified in PipeReg_3
    xxxxxx_yyyyy_00011_00001_xxx...x    // identified in PipeReg_2
\end{verbatim}

\begin{verbatim}
    xxxxxx_00001_xxx...x                // identified in PipeReg_4
    xxxxxx_00011_xxx...x                // identified in PipeReg_3
    xxxxxx_yyyyy_00001_00011_xxx...x    // identified in PipeReg_2
\end{verbatim}

\placedrawing[htbp]{figures/05b_03_forwarding.lp}{The change used to reduce hazards imposed by data dependencies.}{brDelRed}

The structure of the processor is modified as shown in Figure \ref{brDelRed}, where a new level in the pipeline is added and the operands at the ALU inputs are selected using multiplexers with additional inputs. The selection for the augmented multiplexers is computed combinationally as {\tt opSel[3:0]} in {\tt DECODE} and synchronized in {\bf PipeReg2} as {\tt opSel2[3:0]} to avoid an additional delay in selecting ALU operands. Thus, {\tt result} is used as an operand in the EXECUTE stage of the pipeline, if the forwarding mechanism decides, as {\tt result3} connected to inputs 1 or as {\tt result4} to inputs 2 of the selection multiplexers at the ALU operand inputs. The value of {\tt result3} is forwarded when the dependency is identified with the result of the instruction issued one clock cycle before the current one (EX hazard). The value of {\tt result4} is forwarded when the dependency is identified with the result of the instruction issued two clock cycles before the current one (MEM hazard).

\paragraph{Forwarding Truth Table and Priority Logic}
Table \ref{tab:forwarding_logic} summarizes the multiplexer selection codes and hazard conditions for both the left ($rs1$) and right ($rs2$) ALU operands. 

When a register dependency occurs simultaneously across multiple preceding stages---for example, if both \texttt{PipeReg\_3} and \texttt{PipeReg\_4} write to the same register required by \texttt{PipeReg\_2}---the forwarding logic \textbf{must give priority} to \texttt{PipeReg\_3}. Forwarding from \texttt{PipeReg\_4} is only activated if no hazard is detected from \texttt{PipeReg\_3}. This priority rule ensures that the ALU always consumes the most recently generated value in program order.

\begin{table}[H]
  \begin{center}
    \caption{ALU Forwarding Logic and Multiplexer Selection Codes.}
    \label{tab:forwarding_logic}
    \small
    \begin{tabular}{|c|c|l|c|l|}
      \hline
      \textbf{Operand} & \textbf{Selection} & \textbf{Hazard Condition} & \textbf{Multiplexer} & \textbf{Source Selected} \\ 
                       & \textbf{Code}      &                           & \textbf{Input}       &                          \\ \hline \hline
      Left ALU  & \texttt{opSel[3:2] = 00} & No hazard                                            & Input 0 & Register File (\texttt{rf[rs1]}) \\ \cline{2-5}
      Operand   & \texttt{opSel[3:2] = 01} & \texttt{RegWrite\_3} $\land$ (\texttt{rd\_3 == rs1})  & Input 1 & \texttt{result3} (\texttt{PipeReg\_3}) \\ \cline{2-5}
      ($rs1$)   & \texttt{opSel[3:2] = 10} & \texttt{RegWrite\_4} $\land$ (\texttt{rd\_4 == rs1})  & Input 2 & \texttt{result4} (\texttt{PipeReg\_4}) \\ 
                &                          & $\land\ \neg$(\texttt{RegWrite\_3} $\land$ \texttt{rd\_3 == rs1}) & & \\ \hline \hline
      Right ALU & \texttt{opSel[1:0] = 00} & No hazard                                            & Input 0 & Register File (\texttt{rf[rs2]}) \\ \cline{2-5}
      Operand   & \texttt{opSel[1:0] = 01} & \texttt{RegWrite\_3} $\land$ (\texttt{rd\_3 == rs2})  & Input 1 & \texttt{result3} (\texttt{PipeReg\_3}) \\ \cline{2-5}
      ($rs2$)   & \texttt{opSel[1:0] = 10} & \texttt{RegWrite\_4} $\land$ (\texttt{rd\_4 == rs2})  & Input 2 & \texttt{result4} (\texttt{PipeReg\_4}) \\ 
                &                          & $\land\ \neg$(\texttt{RegWrite\_3} $\land$ \texttt{rd\_3 == rs2}) & & \\ \hline
    \end{tabular}
  \end{center}
\end{table}

Formally, the Boolean equations for generating the forwarding multiplexer controls for the left operand ($rs1$) are:
\begin{align}
\text{ForwardA}_{1} &= \texttt{RegWrite\_3} \land (\texttt{rd\_3 == rs1\_2}) \\
\text{ForwardA}_{2} &= \texttt{RegWrite\_4} \land (\texttt{rd\_4 == rs1\_2}) \land \neg \text{ForwardA}_{1}
\end{align}
and symmetrically for the right operand ($rs2$):
\begin{align}
\text{ForwardB}_{1} &= \texttt{RegWrite\_3} \land (\texttt{rd\_3 == rs2\_2}) \\
\text{ForwardB}_{2} &= \texttt{RegWrite\_4} \land (\texttt{rd\_4 == rs2\_2}) \land \neg \text{ForwardB}_{1}
\end{align}
In architectures where register 0 is hardwired to zero (such as RISC-V or MIPS), each hazard condition additionally verifies that $\texttt{rd} \neq 0$ so that writes targeting the constant zero register do not trigger spurious forwarding.

\begin{exemplu}\label{dataDepSimulation}
Data dependency solved by using forwarding is exemplified starting from Example \ref{dataDepSim}:
\begin{verbatim}
            NOP         ;
            VAL(0,1)    ;
            VAL(1,2)    ;
            VAL(2,3)    ;
            VAL(3,4)    ;
            VAL(4,5)    ;
            ADD(5,4,3)  ;
            HALT        ;
            HALT        ;
\end{verbatim}
With forwarding enabled, the processor behaves correctly even without the additional {\tt NOP}s, as follows:
\begin{verbatim}
pc=   0  RF=[x, x, x, x, x, x, x, x]
pc=   1  RF=[x, x, x, x, x, x, x, x]
pc=   2  RF=[x, x, x, x, x, x, x, x]
pc=   3  RF=[x, x, x, x, x, x, x, x]
pc=   4  RF=[x, x, x, x, x, x, x, x]
pc=   5  RF=[x, x, x, x, x, x, x, x]
pc=   6  RF=[1, x, x, x, x, x, x, x]
pc=   7  RF=[1, 2, x, x, x, x, x, x]
pc=   8  RF=[1, 2, 3, x, x, x, x, x]
pc=   8  RF=[1, 2, 3, 4, x, x, x, x]
pc=   8  RF=[1, 2, 3, 4, 5, x, x, x]
pc=   8  RF=[1, 2, 3, 4, 5, 9, x, x]
\end{verbatim}
because the operands for the arithmetic operations are now forwarded  using the two multiplexors to the ALU's inputs.

$\diamond$
\end{exemplu}


\subsection{Control Dependency\index{dependencies!control dependency in pipelined processors}}

The pipelined processor always fetches the instruction immediately after any taken branch.

\begin{exemplu}\label{controlDep}
Let be the following sequence of instructions which contains a unconditioned (relative) jump:
\begin{verbatim}
    i:          XXX     ;
    i+1:        RJMP(12);
    i+2: LB(2)  YYY     ;
    i+3:        UUU     ;
    i+4:        VVV     ;
    ...         ...
    i+j: LB(12) ZZZ     ;
\end{verbatim}
Its execution is supposed to be:
\begin{verbatim}
    XXX     ;
    RJMP(12);
    ZZZ     ;
\end{verbatim}
but in our processor it will be:
\begin{verbatim}
    XXX     ;
    RJMP(12);
    YYY     ;
    ZZZ     ;
\end{verbatim}
because the instruction {\tt YYY} are inserted into the pipe before the execution of the jump instruction completed in {\tt t=i+2} when the jump address is available  (see Table \ref{contrDep}).

\begin{table}[H]
  \begin{center}
    \caption{Hazard generated by the control dependency}
    \label{contrDep}
    {\small
    {\tt
    \begin{tabular}{|c||c|c|c|c|c|c|c|c|}
      \hline
     Instr $\setminus$ Time & t=i  & t=i+1 & t=i+2 & t=i+3 & t=i+4 & t=i+5 & t=i+6 & t=i+7  \\ \hline \hline
        IF                  & XXX0 & RJMP0 & YYY0  & ZZZ0  &       &       &       &        \\ \hline
        OF                  &      & XXX1  & RJMP1 & YYY1  & ZZZ1  &       &       &        \\ \hline
        EX                  &      &       & XXX2  & ---- & YYY2   & ZZZ2  &       &        \\ \hline
        DM                  &       &       &      & XXX3  & ----  & YYY3  & ZZZ3  &        \\ \hline
        WB                  &      &       &       &       & XXX4  & ----  & YYY4  & ZZZ4   \\ \hline
    \end{tabular}
    }
    }
  \end{center}
\end{table}

$\diamond$
\end{exemplu}


\begin{exemplu}
In the following example the effect of interrupt is illustrated together with  the effect of the control dependency due to the unwanted  execution of {\tt VAL(4,33)} positioned after the {\tt RET{30}} instruction.
\begin{verbatim}
            VAL(31,13)      ;
            VAL(2,23)       ;
            VAL(0,13)       ;
            VAL(1,13)       ;
            EI              ;
            ADDV(0,0,1)     ;
            VAL(1,1)        ;
            VAL(2,222)      ;
            NOP             ;
            ADDV(0,0,4)     ;
            HALT            ;
            HALT            ;
            NOP             ;
// subroutine triggered by interrupt
            ADDV(30,30,-1)  ;
            NOP             ;
            NOP             ;
            NOP             ;
            RET(30)         ;
            VAL(4,33)       ;
\end{verbatim}
The result of simulation is:
{\small
\begin{verbatim}
pc=  0  RF=[ x,  x,   x, x,  x, x, x,  x] intState=000 inta=0
pc=  1  RF=[ x,  x,   x, x,  x, x, x,  x] intState=000 inta=0
pc=  2  RF=[ x,  x,   x, x,  x, x, x,  x] intState=000 inta=0
pc=  3  RF=[ x,  x,   x, x,  x, x, x,  x] intState=000 inta=0
pc=  4  RF=[ x,  x,   x, x,  x, x, x, 13] intState=000 inta=0
pc=  5  RF=[ x,  x,  23, x,  x, x, x, 13] intState=000 inta=0
pc=  6  RF=[13,  x,  23, x,  x, x, x, 13] intState=001 inta=0
pc=  7  RF=[13, 13,  23, x,  x, x, x, 13] intState=010 inta=1
pc= 13  RF=[13, 13,  23, x,  x, x, 6, 13] intState=011 inta=0
pc= 13  RF=[13, 13,  23, x,  x, x, 6, 13] intState=100 inta=0
pc= 13  RF=[13, 13,  23, x,  x, x, 6, 13] intState=000 inta=0
pc= 14  RF=[13, 13,  23, x,  x, x, 6, 13] intState=000 inta=0
pc= 15  RF=[13, 13,  23, x,  x, x, 5, 13] intState=000 inta=0
pc= 16  RF=[13, 13,  23, x,  x, x, 5, 13] intState=000 inta=0
pc= 17  RF=[13, 13,  23, x,  x, x, 5, 13] intState=000 inta=0
pc= 18  RF=[13, 13,  23, x,  x, x, 5, 13] intState=000 inta=0
pc=  5  RF=[13, 13,  23, x,  x, x, 5, 13] intState=000 inta=0
pc=  6  RF=[13, 13,  23, x,  x, x, 5, 13] intState=000 inta=0
pc=  7  RF=[13, 13,  23, x,  x, x, 5, 13] intState=000 inta=0
pc=  8  RF=[13, 13,  23, x, 33, x, 5, 13] intState=000 inta=0
pc=  9  RF=[14, 13,  23, x, 33, x, 5, 13] intState=000 inta=0
pc= 10  RF=[14,  1,  23, x, 33, x, 5, 13] intState=000 inta=0
pc= 11  RF=[14,  1, 222, x, 33, x, 5, 13] intState=000 inta=0
pc= 11  RF=[14,  1, 222, x, 33, x, 5, 13] intState=000 inta=0
pc= 11  RF=[18,  1, 222, x, 33, x, 5, 13] intState=000 inta=0
pc= 11  RF=[18,  1, 222, x, 33, x, 5, 13] intState=000 inta=0
pc= 11  RF=[18,  1, 222, x, 33, x, 5, 13] intState=000 inta=0
\end{verbatim}
}
The execution of {\tt VAL(4,33)} is obvious at {\tt t=43}. The program syntax does not assume this.

$\diamond$
\end{exemplu}

There are few solutions for solving the previously emphasized hazard.


\begin{comment}
\subsubsection{Reducing the Branch Delay}

Instead of deciding in the EXECUTION stage of the pipeline if the branch condition is fulfilled, the decision can be moved in OPS FETCH stage. In Figure \ref{brDelRed} the modified structure is presented. The decision is done by testing if the left operand fetched form the register file is zero or not. The adder and the multiplexor used to compute the jump address are moved one stage up.

Why didn't I opt for this solution from the beginning? Because there is a price for the solution we propose. In the initial structure, the zero value could come from the current operation, which could be a logical or arithmetic function. In the new solution we only test the fact that a variable is zero.

Control dependency now produces a reduced hazard. In Table \ref{contrDep1} we notice that the VVV instruction is no longer inserted parasitically in the pipeline. Now instead of flushing three instruction from the pipe when the branch is taken, the system will flush only two instructions.


\begin{table}[H]
  \begin{center}
    \caption{Reducing the hazard generated by the control dependency}
    \label{contrDep1}
    {\small
    {\tt
    \begin{tabular}{|c||c|c|c|c|c|c|c|c|}
      \hline
     Instr $\setminus$ Time & t=i  & t=i+1 & t=i+2 & t=i+3 & t=i+4 & t=i+5 & t=i+6  & t=i+7   \\ \hline \hline
        IF                  & XXX0 & RJMP0 & YYY0  & UUU0  & ZZZ0  &       &        &        \\ \hline
        OF                  &      & XXX1  & RJMP1 & YYY1  & UUU1  & ZZZ1  &        &          \\ \hline
        EX                  &      &       & XXX2  & RJMP2 & YYY2  & UUU2  & ZZZ2   &          \\ \hline
        WB                  &      &       &       & XXX3  & RJMP3 & YYY3  & UUU3   & ZZZ3  \\ \hline
    \end{tabular}
    }
    }
  \end{center}
\end{table}

%\placedrawing[H]{F59.lp}{The  change used to reduce hazards imposed by data and control dependencies.}{brDelRed}


A possible structural disadvantage of this structure may be due to the additional delay introduced by the circuit that detects the zero value at the left output of the register file.

\end{comment}


\subsubsection{Stalling}

By introducing a stall using a {\tt NOP} instruction the sequence of instruction becomes:
\begin{verbatim}
    i:          XXX     ;
    i+1:        RJMP(12);
    i+2:        NOP     ;
    i+3: LB(2)  YYY     ;
    ...         ...
    i+j: LB(12) ZZZ     ;
\end{verbatim}
the execution will be:
\begin{verbatim}
    XXX     ;
    RJMP(12);
    NOP     ;
    ZZZ     ;
\end{verbatim}
This solution results in a branch penalty (the number of stalls introduced during the branch operations in the pipelined processor) of 1 cycle (the  NOP introduced after {\tt RJMP(12)}).

\subsubsection{Reordering}

Sometimes the branch penalty can be reduced by reordering.

Because, in the previous example, the jump is unconditional, which means the execution of the instruction {\tt XXX} does not have implications on the jump instruction, the sequence of instruction can be reordered, as follows:
\begin{verbatim}
    i:          RJMP(12);
    i+1:        XXX     ;
    i+2: LB(2)  YYY     ;
    ...         ...
    i+j: LB(12) ZZZ     ;
\end{verbatim}
thus,  executing the instruction {\tt XXX} after jump no penalty is introduced.

\begin{exemplu}
Reordering is exemplified by the following code:
\begin{verbatim}
            NOP         ;
            VAL(0,-3)   ;
            NOP         ;
            NOP         ;
    LB(1);  NOP         ;
            NOP         ;
            BRNZ(0,1)   ;
            ADDV(0,0,1) ;
            HALT        ;
            HALT        ;
\end{verbatim}
where the increment of the register 0 is specified out of loop, but is executed associated with each branch.
The result of simulation is:
{\small
\begin{verbatim}
pc= 0  RF=[         x, x, x, x, x, x, x, x]
pc= 1  RF=[         x, x, x, x, x, x, x, x]
pc= 2  RF=[         x, x, x, x, x, x, x, x]
pc= 3  RF=[         x, x, x, x, x, x, x, x]
pc= 4  RF=[         x, x, x, x, x, x, x, x]
pc= 5  RF=[4294967293, x, x, x, x, x, x, x] 	
pc= 6  RF=[4294967293, x, x, x, x, x, x, x] 	
pc= 7  RF=[4294967293, x, x, x, x, x, x, x] 	
pc= 4  RF=[4294967293, x, x, x, x, x, x, x] 	
pc= 5  RF=[4294967293, x, x, x, x, x, x, x] 	
pc= 6  RF=[4294967293, x, x, x, x, x, x, x] 	
pc= 7  RF=[4294967294, x, x, x, x, x, x, x] 	
pc= 4  RF=[4294967294, x, x, x, x, x, x, x] 	
pc= 5  RF=[4294967294, x, x, x, x, x, x, x] 	
pc= 6  RF=[4294967294, x, x, x, x, x, x, x] 	
pc= 7  RF=[4294967295, x, x, x, x, x, x, x] 	
pc= 4  RF=[4294967295, x, x, x, x, x, x, x] 	
pc= 5  RF=[4294967295, x, x, x, x, x, x, x] 	
pc= 6  RF=[4294967295, x, x, x, x, x, x, x] 	
pc= 7  RF=[         0, x, x, x, x, x, x, x]
pc= 8  RF=[         0, x, x, x, x, x, x, x]
pc= 9  RF=[         0, x, x, x, x, x, x, x]
pc= 9  RF=[         0, x, x, x, x, x, x, x]
pc= 9  RF=[         1, x, x, x, x, x, x, x]
pc= 9  RF=[         1, x, x, x, x, x, x, x]
pc= 9  RF=[         1, x, x, x, x, x, x, x]
\end{verbatim}
}
Note: 4294967293 = -3.

$\diamond$
\end{exemplu}


\subsubsection{Static Branch Prediction\index{branch prediction!static branch prediction}}

Let us see now how the conditioned jumps, the branches, are managed to be performed efficiently. There are two cases: backward branches or forward branches. Static prediction\index{branch prediction!static branch prediction}presumes that backward branches will be taken and that forward branches will not.

In this case, most of the predictions will be correct. Indeed, if we have a jump back to execute a loop that repeats $n$ times, then once in $n+1$ situations the prediction will have to be corrected.

\begin{exemplu}
A backward branch is one that has a target address that is lower than its own address. For example,
it  refers to the following type of loop:
\begin{verbatim}
    i:          VAL(1,23)   ;
    i+1:        VAL(2,1)    ;
    i+2:  LB(3) SUB(1,1,2)  ;
    i+3:        NZJMP(1,3)  ;
    i+4:        NOP         ;
    i+5:        NOP         ;
    i+6:        ...
\end{verbatim}
This technique can help with prediction accuracy of loops, which are usually backward-pointing branches, and are taken more often than not taken.

$\diamond$
\end{exemplu}

\begin{exemplu}
A forward branch is one that has a target address that is higher than its own address. For example,
it  refers to the following type of loop:
\begin{verbatim}
    i:          VAL(1,23)   ;
    i+1:        VAL(2,1)    ;
    i+2:        SUB(1,1,2)  ;
    i+3:        ZJMP(1,2)   ;
    i+4:        XXX         ;
                ...
    i+j: LB(2)  YYY         ;
\end{verbatim}


$\diamond$
\end{exemplu}

In static prediction, all decisions are made at compile time, before the execution of the program

{\bf Note}: In order to reduce the branch penalty, in the case of unconditional jumps, certain changes can be made in the hardware. For example, the calculation of unconditional jump addresses can be moved to the OP FETCH (OF) level. The execution latency of unconditional jumps is thus reduced by one unit.

\subsubsection{Dynamic Branch Prediction}

Dynamic branch prediction\index{branch prediction!dynamic branch prediction} predicts branch directions and targets based on execution history collected at run-time using dedicated hardware structures. Foundational studies on dynamic branch prediction, saturating up/down counters, and Branch Target Buffers (BTBs)\index{branch target buffer (BTB)} were introduced by J.~E.~Smith \citep{smith1981study} and Lee and A.~J.~Smith \citep{lee1984branch}. We first examine the two classic bimodal counter mechanisms before discussing modern two-level adaptive predictors.

\paragraph{Last-time, one-bit predictor}
is the simplest solution  which uses a two state automaton whose behavior is described in Figure \ref{bp1}. We will code the state {\tt prediction not taken} with 0, and the state {\tt prediction taken} with 1. The automaton is a {\em saturated} 1-bit counter. If is incremented from 0 goes to 1, if is decremented from 1 goes to 0. But if it is incremented from 1, stays in 1, while decremented from 0 stays in 0. This 2-state automaton says whether the branch was recently taken or not. Based on this, the processor fetches the next instruction from the target address or sequential address. If the prediction is wrong, flushes the pipeline and also flips prediction. So, every time a wrong prediction is made, the prediction bit is flipped.

\placedrawing[H]{figures/05b_04_1bit_branch_predictor.lp}{The saturated one-bit up/down counter as last-time branch predictor.}{bp1}


This mechanism always mispredicts the last iteration and the first iteration of a loop branch, thus the accuracy for a loop with $N$ iterations is $(100 \times (N-2)/N)\%$. For large values of $N$ the accuracy of the prediction increases, while for small values this mechanism is not very efficient.


\paragraph{Two-Bit Counter Based Predictor} is a four-state finite automaton. Its behavior is described in Figure \ref{2bp}. The advantage of the two-bit predictor over a one-bit predictor is that a conditional jump has to deviate twice from what it has done most in the past before the prediction changes. For example, a loop-closing conditional jump is mispredicted once rather than twice.

Each state of the automaton is interpreted as follows:
\begin{itemize}
    \item 00: strongly not taken
    \item 01: weakly not taken
    \item 10: weakly taken
    \item 11: strongly taken
\end{itemize}

\placedrawing[H]{figures/05b_05_2bit_branch_predictor.lp}{Two-bit saturated counter based predictor.}{2bp}

The automaton is a {\em saturated} 2-bit counter. It increments for {\tt actually taken} and decrements for {\tt actually not taken}. This predictor changes prediction only on two successive mispredictions.

\paragraph{Step-by-Step Comparison: 1-Bit vs. 2-Bit Predictor on a Loop}
To concretely visualize why a two-bit counter outperforms a one-bit predictor, consider a loop branch that executes 4 iterations ($N = 4$) per invocation before exiting. The dynamic execution sequence of branch outcomes is therefore Taken (T), Taken (T), Taken (T), and Not Taken (NT).

Table \ref{tab:branch_trace} traces two consecutive outer-loop invocations of this loop branch. We assume both predictors start in the Taken state after initial warmup (State 1 for the 1-bit predictor, State 11 / Strongly Taken for the 2-bit counter).

\begin{table}[H]
  \begin{center}
    \caption{Step-by-Step Trace of 1-Bit vs. 2-Bit Dynamic Branch Predictors on a 4-Iteration Loop.}
    \label{tab:branch_trace}
    \small
    \begin{tabular}{|c|c||c|c|c||c|c|c|}
      \hline
      \textbf{Pass \&} & \textbf{Actual} & \multicolumn{3}{c||}{\textbf{1-Bit Predictor}} & \multicolumn{3}{c|}{\textbf{2-Bit Predictor}} \\ \cline{3-8}
      \textbf{Iteration} & \textbf{Outcome} & \textbf{State} & \textbf{Pred} & \textbf{Correct?} & \textbf{State} & \textbf{Pred} & \textbf{Correct?} \\ \hline \hline
      \multicolumn{8}{|l|}{\textit{First Loop Invocation:}} \\ \hline
      Pass 1, Iter 1 & T  & 1 (T)  & T  & \checkmark (Yes) & 11 (ST) & T  & \checkmark (Yes) \\ \hline
      Pass 1, Iter 2 & T  & 1 (T)  & T  & \checkmark (Yes) & 11 (ST) & T  & \checkmark (Yes) \\ \hline
      Pass 1, Iter 3 & T  & 1 (T)  & T  & \checkmark (Yes) & 11 (ST) & T  & \checkmark (Yes) \\ \hline
      Pass 1, Iter 4 & NT & 1 (T)  & T  & $\times$ \textbf{No (Mispred)} & 11 (ST) & T  & $\times$ \textbf{No (Mispred)} \\ \hline \hline
      \multicolumn{8}{|l|}{\textit{Second Loop Invocation (Re-entering Loop):}} \\ \hline
      Pass 2, Iter 1 & T  & 0 (NT) & NT & $\times$ \textbf{No (Mispred)} & 10 (WT) & T  & \checkmark \textbf{Yes (Correct!)} \\ \hline
      Pass 2, Iter 2 & T  & 1 (T)  & T  & \checkmark (Yes) & 11 (ST) & T  & \checkmark (Yes) \\ \hline
      Pass 2, Iter 3 & T  & 1 (T)  & T  & \checkmark (Yes) & 11 (ST) & T  & \checkmark (Yes) \\ \hline
      Pass 2, Iter 4 & NT & 1 (T)  & T  & $\times$ \textbf{No (Mispred)} & 11 (ST) & T  & $\times$ \textbf{No (Mispred)} \\ \hline
    \end{tabular}
  \end{center}
\end{table}

As shown in Table \ref{tab:branch_trace}, the 1-bit predictor incurs \textbf{two mispredictions} per loop invocation:
\begin{enumerate}
    \item When the loop terminates on Iteration 4 (predicts T, actually NT), causing the state bit to flip from 1 to 0.
    \item When the loop is re-entered on Pass 2, Iteration 1 (predicts NT based on the previous loop exit, but the branch is actually T), causing the state bit to flip back to 1.
\end{enumerate}
Consequently, for an $N$-iteration loop, the 1-bit predictor correctly predicts only $N-2$ branch occurrences, giving an accuracy of $(N-2)/N$. For $N = 4$, its steady-state accuracy is only $2/4 = 50\%$.

In contrast, the 2-bit counter predictor incurs only \textbf{a single misprediction} upon loop exit. When the loop terminates on Iteration 4, the counter decrements from Strongly Taken (11) to Weakly Taken (10). Because the most significant bit remains 1, the predictor continues to predict Taken! When the loop re-executes on Pass 2, Iteration 1 is correctly predicted as Taken, and the counter immediately increments back to Strongly Taken (11). Its accuracy is $(N-1)/N$, which is $3/4 = 75\%$ for $N = 4$, and rises to $90\%$ for $N = 10$.

\paragraph{Two-Level Adaptive Predictors and Correlated Branches (Yeh and Patt)}
\label{sec:yeh_patt}
\index{branch prediction!two-level adaptive}
\index{Patt, Yale N.}
\index{Yeh, Tse-Yu}
\index{branch history register (BHR)}
\index{pattern history table (PHT)}
\index{branch correlation}
\index{gshare predictor}
\index{tournament predictor}

While 1-bit and 2-bit bimodal counters track the history of individual branches in isolation, real programs contain branches whose outcomes strongly correlate either with their own recent past outcomes or with the outcomes of other preceding branches (known as \emph{branch correlation}\index{branch correlation}). For example, consider an alternating loop condition or a sequence of dependent conditional statements such as:
\begin{verbatim}
    if (aa == 2)
        aa = 0;
    if (bb == 2)
        bb = 0;
    if (aa != bb) { ... }
\end{verbatim}
If both preceding conditions are true, the third branch is guaranteed not to be taken. An isolated 2-bit counter indexed solely by the branch address cannot capture this inter-branch correlation because it does not record which path led to the branch.

In landmark research, Tse-Yu Yeh and Yale N.~Patt \citep{yeh1991twolevel, yeh1992alternative} solved this limitation by introducing \textbf{Two-Level Adaptive Branch Prediction}. Yale Patt's research group had previously demonstrated with the High Performance Substrate (HPS) architecture \citep{patt1985hps} that aggressive out-of-order execution and wide superscalar dispatch cannot succeed without near-perfect branch prediction, because branch misprediction penalties flush large instruction windows.

The key insight of Yeh and Patt was to decouple dynamic branch prediction into two distinct levels of history:
\begin{enumerate}
    \item \textbf{First Level: Branch History Register (BHR / GHR)}: A $k$-bit shift register that records the actual outcomes of the last $k$ branches (shifting in a 1 for Taken and a 0 for Not Taken). This history can be tracked globally across all branches in a Global History Register (GHR) to capture inter-branch correlation, or locally per branch instruction (BHR) to capture recurring loop patterns.
    \item \textbf{Second Level: Pattern History Table (PHT)}: An array of $2^k$ 2-bit saturating counters. Instead of indexing counters directly by the branch address, the predictor uses the $k$-bit branch history pattern as an index (or combines the branch address with the history pattern) to select a dedicated 2-bit saturating counter.
\end{enumerate}

Yeh and Patt classified two-level adaptive predictors into nine variations based on whether the history register is Global ($G$), Per-branch ($P$), or Set-associative ($S$), and whether the pattern table is Global ($g$), Per-branch ($p$), or Set-associative ($s$)---yielding configurations such as \emph{GAg}, \emph{GAp}, \emph{PAg}, and \emph{PAp} \citep{yeh1992alternative}. By maintaining separate saturating counters for each specific sequence of recent branch outcomes (e.g., distinguishing the behavior after pattern \texttt{1010} from pattern \texttt{1111}), Yeh and Patt demonstrated that two-level adaptive predictors shattered the accuracy ceiling of isolated 2-bit counters, elevating prediction accuracy from $\sim 85\text{--}90\%$ to over $97\%$.

Building upon Yeh and Patt's two-level framework, Scott McFarling \citep{mcfarling1993combining} introduced the \textbf{gshare predictor}, which computes a bitwise XOR of the global branch history register with the branch instruction address (PC) to index the PHT, significantly reducing aliasing (conflict interference) between different branches. McFarling also pioneered \textbf{combining (tournament) predictors}, which use a secondary array of 2-bit counters to dynamically select between two different prediction strategies (such as a local bimodal predictor and a global two-level predictor) based on which predictor has historically demonstrated greater accuracy for that specific branch. Yeh and Patt's two-level adaptive techniques and their derivatives (including gshare, tournament predictors, and modern neural perceptron predictors) formed the core branch prediction engine in virtually all high-performance microprocessors, including the Intel Pentium Pro/II/III/4/Core, AMD K6/Athlon/Zen, and DEC Alpha 21264.


\section{Programming Pipelined Processors}

\begin{exemplu}
The assembly program, written for the generic version of toyRISC, changes the sign of the number stored in {\tt rf[2]} provides in {\tt rf[3]} -7 and then in {\tt rf[1]} back the value of 7.

\begin{verbatim}
            VAL(2, 7);
            CHSGN(3,2);
            CHSGN(1,3);
            HALT;
\end{verbatim}

The perform the same operation in the pipelined version without the forwarding mechanism, the program must take into account the data dependency, is:
\begin{verbatim}
            VAL(2, 7);
            NOP;
            NOP;
            CHSGN(3,2);
            NOP;
            NOP;
            CHSGN(1,3);
            HALT;
            HALT;
\end{verbatim}

Adding the forwarding mechanism, the program becomes:
\begin{verbatim}
            VAL(2, 7);
            CHSGN(3,2);
            CHSGN(1,3);
            HALT;
            HALT;
\end{verbatim}

$\diamond$
\end{exemplu}

\begin{exemplu}
The program which establish by {\tt rf[3] = 1} that the content of {\tt rf[0]} is divisible by 4, else {\tt rf[3] = 0}, is:

\begin{verbatim}
            VAL(0,8);
            VAL(1,3);
            AND(2,0,1);
            NOP;
            NOP;
            BRZ(2,22);
            NOP;
            RJMP(33);
            VAL(3,0);
    LB(22); VAL(3,1);
    LB(33); HALT;
            HALT;
\end{verbatim}
If it runs  preceded by:
\begin{verbatim}
            VAL(0,9);
\end{verbatim}
then the result is {\tt rf[3] = 0}, else if it is preceded by:
\begin{verbatim}
            VAL(0,8);
\end{verbatim}
then the result is {\tt rf[3] = 1}.

The first two {\tt NOP}s in the program are introduced because the forwarding mechanism is implemented only for the operands of the ALU, not for testing the value at the output of the register file by the circuit {\tt zero}. The third {\tt NOP} is due to the control cependency.

$\diamond$
\end{exemplu}

\begin{exemplu}
The program which compute in {\tt rf[4]} the mean values stored in {\tt rf[0], rf[1], rf[2], rf[3]} is:
\begin{verbatim}
            ADD(4,0,1);
            ADD(4,4,2);
            ADD(4,4,3);
            ASH(4,4);
            ASH(4,4);
            HALT;
            HALT;
\end{verbatim}
If it runs  preceded by:
\begin{verbatim}
            VAL(0,10);
            VAL(1,12);
            VAL(2,3);
            VAL(3,44);
\end{verbatim}
then the result if {\tt rf[4] = 17}.

$\diamond$
\end{exemplu}

\section{Superscalar Processor\index{superscallar processor: a processor  designed with several execution units}}

If the processor is designed with several execution units (for example: one integer ALU, two floating-point units, two load/store units), the potential parallelism thus obtained should be activated at the highest possible level. This means that in each clock cycle a maximum number of resources should be active, if possible all of them.

In this case, the processor, called  {\em superscalar}, can execute instructions in an order imposed by the availability of data and execution units, rather than by their initial sequence in the  program. For example, the {\em PowerPC 970} processor  fetches and decodes up to eight instructions, dispatch up to five to reserve stations (register files), issue up to eight to the execution units and retire up to five per cycle.

\subsection{Register renaming\index{renaming register: technique applied when sequential execution is no longer mandatory}}

The sequence in which the instructions are executed must be strictly respected when there are data dependencies. But when certain independent sequences can be identified, strictly sequential execution is no longer mandatory, especially in the case of a superscalar processor that allows parallelism at the level of instructions. We can apply in such cases the {\em register renaming} technique.

\begin{exemplu}
Consider this sequence of code running on an out-of-order processor:
\begin{verbatim}
    READ(0,4)   ;   // read in rf[0] from the address in rf[4]
    ADD(0,0,12) ;   // rf[0] <= rf[0] + rf[12]
    STORE(0,7)  ;   // dataMemory[rf[7]] <= rf[0]
    READ(0,5)   ;
    ADD(0,0,6)  ;
    STORE(0,21) ;
\end{verbatim}
The instructions in the last three lines are {\bf independent} of the first three lines. The processor cannot execute {\tt STORE(0,21)} until  {\tt STORE(0,7)} is done.
We can eliminate this restriction involving in our code a supplementary register, as follows:
\begin{verbatim}
    READ(0,4)   ;   // read in rf[0] from the address in rf[4]
    ADD(0,0,12) ;   // rf[0] <= rf[0] + rf[12]
    STORE(0,7)  ;   // dataMemory[rf[7]] <= rf[0]
    READ(2,5)   ;
    ADD(2,2,6)  ;
    STORE(2,21) ;
\end{verbatim}
The first three instructions involves {\tt rf[0]}, while the last three instructions work on the register {\tt rf[2]}, allowing the last three instructions to be executed in parallel with the first three.

$\diamond$
\end{exemplu}


\subsection{Out-of-Order Execution\index{out-of-order execution: instruction execution in an order governed by the data availability}}

Out-of-order execution  is a mechanism used in  high-performance  processing units to make use efficiently  of instruction cycles. Using this mechanism, a processor executes instructions in an order governed by the availability of input data and execution units, rather than by their original order in a program. Thus the processor  avoids being idle while waiting for the preceding instruction to complete and can, in the meantime, proceed to execution of the next instructions that are able to run immediately and independently.
The mechanism is efficient for a processor with multiple execution units {\bf with speed difference between instructions}. It is the case of superscalar processor with ISA performing integer and floating-point arithmetic operations. There are also speed difference between arithmetic and logic instructions and memory access instructions.

While in a pipelined scalar processor an instruction is executed in the following steps:
\begin{enumerate}
\item Instruction fetch.
\item Operands fetch, if input operands are available in processor's register file, else the processor stalls until they are available.
\item The instruction is executed.
\item The results is write back to the register file.
\end{enumerate}
in a superscalar processor the out-of-order execution is requested. It consists of the following steps:
\begin{enumerate}
\item Instruction fetch.
\item Instruction dispatch to an instruction queue (also called instruction buffer or reservation stations).
\item The instruction waits in the queue until its input operands are available. The instruction can leave the queue before older instructions.
\item The instruction is issued to the appropriate idle functional unit.
\item The results are queued.
\item Only after all older instructions have their results written back to the register file, then this result is written back to the register file.
\end{enumerate}

The benefit of out-of-order execution processing grows as the instruction pipeline deepens and the speed difference between main memory or cache memory and the processor widens. On modern machines, the processor runs few times faster than the cache memory and many times faster then system memory, so during the time an in-order processor  waits for data, a large number of instructions could be executed.

\begin{exemplu}\label{warwaw}
In the following sequence of instructions (see \citep{Patterson05}, p. 185):
\begin{verbatim}
    DIV(0,2,4)  ;
    ADD(6,0,8)  ;
    STR(6,1)    ;   // store
    SUB(8,9,7)  ;
    MUL(6,9,8)  ;
\end{verbatim}
here are the following dependencies:
\begin{enumerate}
    \item between ADD and SUB if SUB finishes before ADD starts (a WAR hazard) is an antidependence
    \item between ADD and MUL if ADD finishes later than MUL (WAW hazard) is an output dependence
    \item between DIV and ADD a true data dependency (a RAW hazard)
    \item between SUB and MUL a true data dependency (a RAW hazard)
    \item between ADD and STORE a true data dependency (a RAW hazard)
\end{enumerate}
First, let us apply register renaming technique involving two additional registers, {\tt rf(10)} and {\tt rf(11)}:
\begin{verbatim}
    DIV(0, 2, 4)    ;
    ADD(10,0, 8)    ;
    STR(10,1)       ;
    SUB(11,9, 7)    ;
    MUL(6, 9, 11)   ;
\end{verbatim}
Any WAR and WAW dependencies are now removed statically by the compiler. For RAW dependency the forwarding mechanism works but only in a processor with one integer ALU. What can be done for a superscalar processor with floating point arithmetic?

$\diamond$
\end{exemplu}

The hardware added for out-of-order execution is big sized and complex. With the increase in the number of pipeline levels, the costs (area plus complexity in use) increase. Multiprocessing and multithreading will allow more efficient alternative solutions.

\begin{comment}
\subsubsection{Floating-point representation of real numbers}

To expand the range in which the numbers are represented, the following form is used:
$$N = sgn(1 + frac) \times 2^{exp - 127}$$
where:
$$sgn \in \{-,+\} = \{0,1\}$$
$$frac \in [0, 1)$$
$$exp \in [0,255]$$
The binary representation is:
$$\langle sgn | exp | frac \rangle$$
For example the 32-bit version:
$$0\_10000001\_10000000000000000000000$$
represents the number: $+ (1+0.5) \times 2^{129-127} = 6$

More at: \url{https://www.geeksforgeeks.org/ieee-standard-754-floating-point-numbers/}

For the purpose we are pursuing, it is enough to understand that the operation with represented floating-point numbers involves a pipeline or a sequence of elementary operations of one clock cycle. The number of cycles will be different depending on the operation, minimum for addition and maximum for division.
\end{comment}

\subsubsection{Dynamic Scheduling with the Scoreboard (CDC 6600)}\label{sec:scoreboard}\index{Scoreboard!CDC 6600 dynamic scheduling}\index{Cray, Seymour (1925-1996)}\index{Thornton, James E. (1925-2005)}

The earliest dynamic scheduling technique for out-of-order execution was the \emph{Scoreboard}, developed by Seymour Cray and James E. Thornton for the Control Data Corporation (CDC) 6600 supercomputer in 1964 \citep{thornton1964, thornton1970}. The CDC 6600 incorporated multiple independent functional units (floating-point adders, multipliers, a divider, shift units, and integer ALUs) operating in parallel. 

The scoreboard is a centralized hardware control unit that maintains the status of all active instructions, functional units, and registers. It dynamically steers instruction progress through four execution stages:
\begin{enumerate}
    \item \textbf{Issue}: Decodes the instruction and checks for \textbf{structural hazards} and \textbf{WAW hazards} (output dependencies). If the required functional unit is busy, or if another previously issued instruction has the same destination register ($rd$), the instruction issue stage stalls. No subsequent instructions can issue until the hazard resolves.
    \item \textbf{Read Operands}: Monitors \textbf{RAW hazards} (true data dependencies). An instruction remains in this stage until all of its source operands are available in the register file---meaning no currently active functional unit is computing a value destined for those registers. Once the operands are valid, the functional unit reads them directly from the register file.
    \item \textbf{Execution}: The functional unit performs the requested operation over its required latency (e.g., 1 cycle for integer addition, 4 cycles for floating-point addition, 10 cycles for floating-point multiplication). When execution completes, the unit notifies the scoreboard.
    \item \textbf{Write Result}: Checks for \textbf{WAR hazards} (anti-dependencies). Before a functional unit can write its result back to the register file, the scoreboard checks whether any earlier issued instruction has yet to read its source operand from this destination register. If such an instruction exists, the writeback stage stalls until the dependent instruction has completed its \textit{Read Operands} phase!
\end{enumerate}

Although the scoreboard achieved unprecedented performance for its time by enabling independent operations to complete out of order, it suffered from two key limitations:
\begin{itemize}
    \item \textbf{WAR and WAW Stalls}: Because registers are referenced strictly by their architectural names without renaming, instructions must stall on both WAR and WAW hazards.
    \item \textbf{Centralized Bottleneck}: All status tracking and hazard evaluations occur in a centralized scoreboard, restricting issue bandwidth and limiting scalability to deeply pipelined or wider-issue superscalar cores.
\end{itemize}

\subsubsection{Tomasulo's Algorithm and Dynamic Register Renaming (IBM 360/91)}\label{sec:tomasulo}\index{Tomasulo's algorithm: dynamic scheduling in out-of-order execution}\index{Tomasulo, Robert (1934-2008)}

To overcome the performance bottlenecks of scoreboarding, Robert M. Tomasulo invented an algorithm for dynamic instruction scheduling with register renaming at IBM in 1967 for the System/360 Model 91 floating-point unit \citep{tomasulo1967}. Tomasulo's algorithm introduces three revolutionary microarchitectural concepts:
\begin{enumerate}
    \item \textbf{Distributed Reservation Stations}: Instead of a centralized control structure, small storage buffers called \emph{Reservation Stations} (RS) are placed directly at the input of each functional unit. Instructions are dispatched immediately into an available reservation station, freeing the instruction issue logic to fetch and decode subsequent instructions.
    \item \textbf{Dynamic Register Renaming}: Register specifiers are replaced by \emph{tags} identifying the reservation station that will produce the result. If a source operand is not yet available at issue time, the reservation station stores the tag of the producing unit rather than the operand value. Because operands are buffered and referenced by tags rather than architectural register names, \textbf{WAR and WAW hazards are completely eliminated without stalling}.
    \item \textbf{Common Data Bus (CDB)}: Results are not written point-to-point into the register file. Instead, when a functional unit finishes, it broadcasts its reservation station tag along with the computed data over a shared broadcast bus called the \emph{Common Data Bus} (CDB). All reservation stations waiting for that tag snoop the bus and latch the result simultaneously, bypassing the register file entirely.
\end{enumerate}

Tomasulo's algorithm coordinates execution through three distinct phases:
\begin{enumerate}
    \item \textbf{Issue}: Checks for a structural hazard (is there a free reservation station for this instruction's functional unit type?). If a reservation station is free, the instruction is issued. If a source operand is already available in the register file, its value is copied directly into the reservation station's value field ($V_j$ or $V_k$). If the operand is currently being computed by another instruction, the tag of the producing reservation station is copied into the tag field ($Q_j$ or $Q_k$). Finally, the destination register in the register status table is updated to point to this reservation station ($Q_i$), dynamically renaming the register.
    \item \textbf{Execute}: When all required source operands are available in the reservation station (i.e., $Q_j = 0$ and $Q_k = 0$), the functional unit begins execution. While waiting, the reservation station snoops the CDB on every cycle to capture pending values when their producing tags are broadcast, resolving \textbf{RAW hazards} dynamically.
    \item \textbf{Write Result}: Upon completion, the functional unit broadcasts its result and tag over the CDB. Any reservation station waiting on that tag copies the value into its $V$ field and clears its $Q$ field. If the register status table indicates that this reservation station is still the latest producer for the destination register ($Q_i == \text{this RS}$), the register file is updated. Finally, the reservation station is marked free ($Busy = 0$).
\end{enumerate}

Each entry in a Reservation Station maintains the following status fields \citep{Patterson19}:
\begin{itemize}
    \item \textbf{Op}: The operation to perform on source operands $S_1$ and $S_2$.
    \item \textbf{Qj, Qk}: The reservation station tags that will produce the corresponding source operands; a value of zero indicates that the source operand is already available in $V_j$ or $V_k$, or is unnecessary.
    \item \textbf{Vj, Vk}: The values of the source operands. Note that for each operand, only one of the $V$ field or $Q$ field is valid. For memory store instructions, $V_k$ holds the data value to be stored.
    \item \textbf{A}: Holds memory address calculation information for load and store instructions. Initially, the immediate offset is stored here; after effective address calculation, the memory address is stored here.
    \item \textbf{Busy}: A 1-bit flag indicating whether this reservation station and its functional unit are currently occupied.
\end{itemize}

Complementing the reservation stations, the processor's register file includes a register result status field, \textbf{Qi}:
\begin{itemize}
    \item \textbf{Qi}: The identifier of the reservation station producing the result to be stored into register $R_i$. If $Q_i = 0$ (or blank), no currently active instruction is writing to $R_i$, meaning the value stored in the register file is architecturally valid.
\end{itemize}

\subsubsection{Speculative Execution, Precise Interrupts, and the Reorder Buffer}\label{sec:rob}
\index{reorder buffer (ROB)!speculative execution}\index{precise interrupts!out-of-order execution}\index{Smith, James E.}\index{Pleszkun, Andrew R.}\index{Sohi, Gurindar S.}

While Tomasulo's algorithm dynamically eliminates false dependencies (WAR and WAW) and tolerates long execution latencies, it possesses a fundamental architectural limitation: \textbf{imprecise exceptions}. In classic Tomasulo scheduling, results are written out of order directly to the architectural register file as soon as each functional unit completes. 

If an instruction triggers a hardware exception (such as an arithmetic overflow, divide by zero, or virtual memory page fault), or if a conditional branch is speculatively mispredicted, subsequent instructions that executed out of order may have already permanently overwritten architectural registers. Consequently, the processor cannot reconstruct the exact state of the program at the point of the fault, making software debugging, page-fault handling, and speculative branch recovery impossible.

To resolve this critical flaw, James E. Smith and Andrew R. Pleszkun introduced the \textbf{Reorder Buffer (ROB)} in 1988 \citep{smith1988precise}. Gurindar S. Sohi extended this approach with the Register Update Unit (RUU) \citep{sohi1990instruction}, and Smith and Sohi synthesized the foundational framework for modern speculative superscalar microarchitectures \citep{smith1991}.

The core insight of the Reorder Buffer is to \textbf{decouple execution completion from architectural state update}:
\begin{itemize}
    \item Instructions are dispatched and issued \textbf{in strict program order}.
    \item Instructions execute and produce results \textbf{out of order}.
    \item Completed results are buffered \textbf{speculatively} in the Reorder Buffer.
    \item Instructions commit and update the permanent architectural register file and memory \textbf{strictly in program order}.
\end{itemize}

\paragraph{ROB Organization and Entry Fields}
The Reorder Buffer is organized as a circular FIFO queue managed by head (commit) and tail (issue) pointers. When an instruction is decoded and dispatched, it is assigned a unique sequential slot at the tail of the ROB. The index of this ROB slot acts as the dynamic renaming tag for the instruction's destination register. 

Each entry in the Reorder Buffer contains the following fields:
\begin{itemize}
    \item \textbf{Type}: Identifies the instruction class (e.g., branch, store, integer ALU, floating-point, load).
    \item \textbf{Destination}: Specifies the target architectural register number ($rd$) for register-writing instructions, or the computed memory destination for store instructions.
    \item \textbf{Value}: Holds the speculative data result computed by the functional unit once execution completes.
    \item \textbf{Ready}: A 1-bit status flag indicating whether the functional unit has finished execution and broadcast its result to the ROB.
    \item \textbf{Exception}: Status flags indicating whether an arithmetic overflow, memory alignment fault, or page fault was raised during execution.
\end{itemize}

\paragraph{Four-Stage Speculative Out-of-Order Execution}
With a Reorder Buffer, instruction execution proceeds through four distinct stages:
\begin{enumerate}
    \item \textbf{Issue / Dispatch (In-Order)}: The next instruction in program order is decoded. If both a reservation station and an entry at the tail of the ROB are available, the instruction is allocated both resources. The destination register in the register status table is updated to point to the newly assigned ROB tag ($Q_i = \text{ROB}_{\text{tail}}$). If a source operand is already available in the architectural register file, it is read; if it is currently in flight in an active ROB entry, it is copied from the ROB (if ready) or the producing ROB tag is recorded in the reservation station.
    \item \textbf{Execute (Out-of-Order)}: As soon as all required source operands become available in the reservation station (via CDB forwarding or direct dispatch), the functional unit executes the operation.
    \item \textbf{Write Result (Out-of-Order)}: When the functional unit completes, it broadcasts its computed result and ROB tag over the Common Data Bus (CDB). The result is written into the assigned ROB entry, and its \textit{Ready} bit is set to 1. Simultaneously, any waiting reservation stations snoop the CDB and capture the value. \textbf{The architectural register file is not touched during this stage}.
    \item \textbf{Commit / Retire (Strictly In-Order)}: Only the instruction at the \textbf{head} of the ROB is evaluated for retirement. When the head entry has $\textit{Ready} = 1$:
    \begin{itemize}
        \item \textit{Normal Completion}: If no exception occurred, the speculative value stored in the ROB entry is written permanently into the architectural register file (or committed to data cache/memory if it is a store instruction). The head pointer advances, freeing the ROB entry.
        \item \textit{Branch Misprediction}: If the head instruction is a conditional branch that was mispredicted, all speculative entries currently in the ROB are instantaneously flushed in a single clock cycle, the processor redirects the program counter to the alternate target, and execution restarts without any corrupted architectural registers!
        \item \textit{Precise Exception}: If an exception flag is set, the pipeline and all subsequent ROB entries are discarded. Because no instruction past the faulting instruction has been allowed to commit to the architectural register file, the processor state is 100\% \textbf{precise}, allowing the operating system exception handler to service the interrupt and resume execution correctly.
    \end{itemize}
\end{enumerate}

\subsubsection{Comparative Summary: Scoreboard, Tomasulo, and Reorder Buffer}\label{sec:ooo_comparison}

Table \ref{tab:scoreboard_vs_tomasulo} contrasts the structural and operational trade-offs across the three landmark generations of dynamic instruction scheduling.

\begin{table}[H]
  \begin{center}
    \caption{Architectural Comparison of Dynamic Scheduling Approaches.}
    \label{tab:scoreboard_vs_tomasulo}
    \small
    \begin{tabular}{|l|l|l|l|}
      \hline
      \textbf{Feature} & \textbf{Scoreboard (1964)} & \textbf{Tomasulo (1967)} & \textbf{Tomasulo + ROB (1988)} \\ 
                       & \citep{thornton1964}       & \citep{tomasulo1967}     & \citep{smith1988precise, smith1991} \\ \hline \hline
      \textbf{Control Unit}      & Centralized scoreboard      & Distributed RS           & Distributed RS + Circular ROB \\ \hline
      \textbf{Renaming}          & None (architectural regs)   & Via RS tags \& $Q_i$     & Dynamic via ROB entry tags \\ \hline
      \textbf{RAW Hazards}       & Stalls \textit{Read Op}     & Waits in RS; snoops CDB  & Waits in RS; snoops CDB/ROB \\ \hline
      \textbf{WAR Hazards}       & \textbf{Stalls} Writeback   & \textbf{Eliminated}      & \textbf{Eliminated} by renaming \\ \hline
      \textbf{WAW Hazards}       & \textbf{Stalls} Issue       & \textbf{Eliminated}      & \textbf{Eliminated} by renaming \\ \hline
      \textbf{Speculative State} & In register file only       & In RS / register file    & Buffered in ROB until commit \\ \hline
      \textbf{Commit Order}      & Out-of-order                & Out-of-order             & \textbf{Strictly in-order} \\ \hline
      \textbf{Exceptions}        & \textbf{Imprecise}          & \textbf{Imprecise}       & \textbf{Precise interrupts} \\ \hline
      \textbf{Branch Recovery}   & Stalls at branch            & Stalls at branch         & \textbf{Instant flush of ROB} \\ \hline
      \textbf{Hardware Cost}     & Moderate centralized logic  & High (RS + CDB routing)  & Highest (RS + CDB + FIFO ROB) \\ \hline
    \end{tabular}
  \end{center}
\end{table}

\section{Problems}

\begin{problem}
Write a program which loads 5  prime numbers in the register file and then compute their average value. Do it for the three versions of toyRISC processor, the generic, the pipelined and the pipelined with forward. Test the programs on the three  versions of toyRISC presented and designed in Chapter \ref{toyRiscDesign}.

$\diamond$
\end{problem}

\begin{problem}
Write, for the three versions of toyRISC, a program which increment continuously the content of the register 12 starting from value 22, and accepts an interrupt which uses a subroutine designed to  square the content of the register 12 and to send the result to the external memory at de address stored in the register 13.

$\diamond$
\end{problem}


\begin{problem}
Design, for the three versions of toyRISC,  a program which loads in data memory, starting form the address 0, the the first 32 numbers of the Fibonacci sequence.

$\diamond$
\end{problem}

\begin{problem}
Write, for the three versions of toyRISC,  the program which compute $10!$.

$\diamond$
\end{problem}

\begin{problem}
\textbf{Pipeline Speedup and Imbalance Analysis}:
Consider an unpipelined processor where the major datapath operations have the following latencies:
Instruction Fetch ($\text{IF}$) = $250\,\text{ps}$,
Operand Fetch / Decode ($\text{OF}$) = $150\,\text{ps}$,
Execute ($\text{EX}$) = $320\,\text{ps}$,
Memory Access ($\text{MEM}$) = $380\,\text{ps}$,
Write Back ($\text{WB}$) = $150\,\text{ps}$.
Suppose the processor is partitioned into a 5-stage pipeline ($\text{IF}, \text{OF}, \text{EX}, \text{MEM}, \text{WB}$). Adding pipeline registers between each stage introduces a setup and clock-to-Q latch overhead of $t_{\text{reg}} = 30\,\text{ps}$ per stage.
\begin{enumerate}
    \item Determine the minimum clock period and operating frequency for both the unpipelined and pipelined processors.
    \item Calculate the ideal speedup achieved by pipelining for an arbitrarily large number of instructions ($N \to \infty$), assuming zero hazard stalls.
    \item Suppose that due to data load delays and branch mispredictions, dynamic execution incurs an average penalty of 0.25 stall cycles per instruction. What is the actual speedup achieved over the unpipelined core?
\end{enumerate}

\textbf{Solution}:
\begin{enumerate}
    \item For the unpipelined processor, an instruction executes sequentially through all five functional blocks in a single clock cycle:
    \[
    T_{\text{unpipe}} = 250 + 150 + 320 + 380 + 150 = 1250\,\text{ps} \quad (f_{\text{unpipe}} = 800\,\text{MHz})
    \]
    For the 5-stage pipelined processor, the clock cycle is limited by the critical path of the slowest (bottleneck) stage plus the pipeline register overhead:
    \[
    T_{\text{pipe}} = \max(250, 150, 320, 380, 150) + t_{\text{reg}} = 380\,\text{ps} + 30\,\text{ps} = 410\,\text{ps} \quad (f_{\text{pipe}} \approx 2.439\,\text{GHz})
    \]
    \item Under ideal conditions with no stalls, the pipelined processor retires 1 instruction per cycle ($\text{CPI}_{\text{pipe}} = 1.0$), while the unpipelined processor executes 1 instruction per cycle with clock period $T_{\text{unpipe}}$ ($\text{CPI}_{\text{unpipe}} = 1.0$). The ideal speedup is:
    \[
    S_{\text{ideal}} = \frac{T_{\text{unpipe}}}{T_{\text{pipe}}} = \frac{1250\,\text{ps}}{410\,\text{ps}} \approx 3.049\times
    \]
    Notice that the speedup is significantly less than the number of stages ($k = 5$) due to stage imbalance (the $380\,\text{ps}$ MEM stage vs. $150\,\text{ps}$ stages) and register overhead ($30\,\text{ps}$).
    \item With $0.25$ stall cycles per instruction, the average CPI of the pipelined processor increases to $\text{CPI}_{\text{actual}} = 1.0 + 0.25 = 1.25$. The actual speedup becomes:
    \[
    S_{\text{actual}} = \frac{\text{CPI}_{\text{unpipe}} \times T_{\text{unpipe}}}{\text{CPI}_{\text{actual}} \times T_{\text{pipe}}} = \frac{1.0 \times 1250\,\text{ps}}{1.25 \times 410\,\text{ps}} = \frac{1250}{512.5} \approx 2.439\times
    \]
\end{enumerate}
$\diamond$
\end{problem}

\begin{problem}
\textbf{Forwarding Unit Hazard Equations and Priority}:
In the pipelined toyRISC processor with forwarding, the Execute stage ($\text{EX}$) receives operands selected from the register file, from \texttt{PipeReg\_3} (holding the result of the instruction issued 1 cycle earlier), or from \texttt{PipeReg\_4} (holding the result of the instruction issued 2 cycles earlier).
\begin{enumerate}
    \item Specify the Boolean logic expressions for selecting the left ALU operand multiplexer control signals ($\text{ForwardA}$).
    \item Consider the following assembly code snippet:
\begin{verbatim}
    ADD(1, 2, 3)  ;   // rf[1] <= rf[2] + rf[3]
    ADD(1, 4, 5)  ;   // rf[1] <= rf[4] + rf[5]
    SUB(6, 1, 7)  ;   // rf[6] <= rf[1] - rf[7]
\end{verbatim}
Trace the contents of \texttt{PipeReg\_3}, \texttt{PipeReg\_4}, and \texttt{PipeReg\_2} when \texttt{SUB} enters the Execute stage. Explain why a naive forwarding unit that checks \texttt{PipeReg\_4} without masking by \texttt{PipeReg\_3} would produce an erroneous architectural result.
\end{enumerate}

\textbf{Solution}:
\begin{enumerate}
    \item Let $\text{ForwardA} = 01$ select \texttt{result3} from \texttt{PipeReg\_3}, $\text{ForwardA} = 10$ select \texttt{result4} from \texttt{PipeReg\_4}, and $\text{ForwardA} = 00$ select the register file output. The Boolean conditions are:
    \begin{align*}
    \text{ForwardA}_{01} &= \texttt{RegWrite\_3} \land (\texttt{rd\_3 == rs1\_2}) \\
    \text{ForwardA}_{10} &= \texttt{RegWrite\_4} \land (\texttt{rd\_4 == rs1\_2}) \land \neg(\texttt{RegWrite\_3} \land (\texttt{rd\_3 == rs1\_2}))
    \end{align*}
    \item When \texttt{SUB(6,1,7)} reaches the Execute stage:
    \begin{itemize}
        \item \texttt{PipeReg\_2} holds \texttt{SUB(6,1,7)} with source register $\texttt{rs1} = 1$.
        \item \texttt{PipeReg\_3} holds the second \texttt{ADD(1,4,5)} with destination register $\texttt{rd} = 1$.
        \item \texttt{PipeReg\_4} holds the first \texttt{ADD(1,2,3)} with destination register $\texttt{rd} = 1$.
    \end{itemize}
    Both \texttt{PipeReg\_3} and \texttt{PipeReg\_4} have $\texttt{rd} = 1$, matching $\texttt{rs1} = 1$. However, in program order, the second instruction (\texttt{ADD(1,4,5)}) overwrote the value computed by the first. If the hardware does not prioritize \texttt{PipeReg\_3} over \texttt{PipeReg\_4}, the forwarding unit might select the stale result from the older instruction in \texttt{PipeReg\_4}, violating RAW data integrity. Masking the \texttt{PipeReg\_4} condition with $\neg(\texttt{RegWrite\_3} \land (\texttt{rd\_3 == rs1\_2}))$ guarantees that the newest program value is always consumed.
\end{enumerate}
$\diamond$
\end{problem}

\begin{problem}
\textbf{Branch Predictor Performance and CPI Penalty}:
A pipelined processor executes a dynamic benchmark where 20\% of all executed instructions are conditional branches. Due to the branch resolution latency in the pipeline, each mispredicted branch incurs a 3-cycle flush penalty, whereas a correctly predicted branch incurs 0 penalty cycles.
Assume that 70\% of all executed conditional branches are taken in practice.
\begin{enumerate}
    \item Calculate the effective CPI (assuming ideal $\text{CPI}_{\text{base}} = 1.0$) for a static branch predictor that predicts Always Not Taken.
    \item Calculate the effective CPI for a 1-bit dynamic predictor that achieves an 82\% prediction accuracy on this workload.
    \item Calculate the effective CPI for a 2-bit saturating counter dynamic predictor that achieves a 93\% prediction accuracy.
    \item What is the percentage speedup of the 2-bit predictor over the static Always Not Taken predictor?
\end{enumerate}

\textbf{Solution}:
\begin{enumerate}
    \item The general formula for CPI including branch misprediction penalty is:
    \[
    \text{CPI} = \text{CPI}_{\text{base}} + (\text{Branch Frequency}) \times (\text{Misprediction Rate}) \times (\text{Penalty})
    \]
    For the static Always Not Taken predictor, every branch that is actually taken is mispredicted. Since 70\% of branches are taken, the misprediction rate is $0.70$:
    \[
    \text{CPI}_{\text{static}} = 1.0 + 0.20 \times 0.70 \times 3 = 1.0 + 0.420 = 1.420
    \]
    \item For the 1-bit dynamic predictor with 82\% accuracy, the misprediction rate is $1.0 - 0.82 = 0.18$ (18\%):
    \[
    \text{CPI}_{\text{1-bit}} = 1.0 + 0.20 \times 0.18 \times 3 = 1.0 + 0.108 = 1.108
    \]
    \item For the 2-bit saturating counter predictor with 93\% accuracy, the misprediction rate is $1.0 - 0.93 = 0.07$ (7\%):
    \[
    \text{CPI}_{\text{2-bit}} = 1.0 + 0.20 \times 0.07 \times 3 = 1.0 + 0.042 = 1.042
    \]
    \item The execution speedup of the 2-bit predictor over the static Always Not Taken predictor is:
    \[
    \text{Speedup} = \frac{\text{CPI}_{\text{static}}}{\text{CPI}_{\text{2-bit}}} = \frac{1.420}{1.042} \approx 1.3628 \implies \mathbf{36.3\% \text{ speedup}}
    \]
\end{enumerate}
$\diamond$
\end{problem}


%\placedrawing{IMEC.lp}{}{} 